\documentclass[journal]{IEEEtran}
\usepackage{amsmath,amssymb,amsthm}
\usepackage{booktabs}
\usepackage{cite}
\usepackage[hidelinks]{hyperref}
\newif\ifanon \anonfalse
\IfFileExists{anon.flag}{\anontrue}{}
\hypersetup{pdftitle={Fixed-Past Threshold Reduction for Finite Games with Repetition Rules, with an Exact Instantiation for FIDE Chess},pdfkeywords={repetition rules, graph-history interaction, reachability games, endgame tablebases, certifying algorithms}}
\ifanon\else\hypersetup{pdfauthor={Muaawia Jannat}}\fi
\usepackage{microtype}

\newtheorem{theorem}{Theorem}
\newtheorem{lemma}{Lemma}
\newtheorem{corollary}{Corollary}
\newtheorem{proposition}{Proposition}
\theoremstyle{definition}
\newtheorem{definition}{Definition}
\newtheorem{example}{Example}
\theoremstyle{remark}
\newtheorem{remark}{Remark}
\newcommand{\sat}{\operatorname{sat}}
\newcommand{\N}{\mathbb{N}}

\begin{document}
\title{Fixed-Past Threshold Reduction for Finite Games with Repetition Rules, with an Exact Instantiation for FIDE Chess}
\ifanon
\author{Anonymous authors}
\else
\author{Muaawia~Jannat%
\thanks{M. Jannat is an independent researcher in Calgary, Canada (e-mail: jannat.muaawia@ucalgary.ca; ORCID 0009-0006-4753-3156).}}
\fi
\markboth{IEEE Transactions on Games}{\ifanon Anonymous\else Jannat\fi: Fixed-Past Threshold Reduction for Games with Repetition Rules}
\maketitle

\begin{abstract}
Repetition rules make the value of a game position depend on the history that reached it. We study finite deterministic two-player games in which a position that has recurred may be claimed as a draw, a position that has recurred often enough is drawn automatically, quiet moves advance a bounded clock, and irreversible moves enter exactly valued subgames. For a query with an arbitrary supplied past we prove that two frozen sets of positions, those the past has already brought to within one occurrence of a claimable draw and those within one occurrence of an automatic draw, determine the exact win/draw/loss value, with no occurrence counting during the solve. Under alternating ownership the second set is unnecessary once the query position's own status is checked. The proof identifies the latest winning clock at each position with the attractor rank of a reachability game on positions, which yields a winning strategy that is positional in the position alone and never revisits one, so no future repetition event can arise, including a defender's option to claim before making a repeating move; the same identification gives a bucket algorithm linear in the arena size plus the clock bound and a locally checkable certificate. Machine-checked counterexamples show that neutrality of repetition, monotonicity of the clock and alternation are each necessary. Two independently written implementations agree on every one of 850 million queries over exhaustively enumerated arenas. For FIDE chess the reduction is instantiated by a solver for every pawnless three- and four-piece material under the exact repetition and move-clock rules; its fresh-history values agree with the Syzygy tables on every compared position, and its deadline equals 100 minus the exact distance to zeroing on every winning position that is not an immediate capture or mate. The theorem explains why engine repetition heuristics are value-correct with full history and quantifies how often a history-free probe is wrong.
\end{abstract}

\begin{IEEEkeywords}
repetition rules, graph-history interaction, reachability games, endgame tablebases, certifying algorithms
\end{IEEEkeywords}

\section{Introduction}

A transposition table identifies plays that reach the same position. Repetition rules make those plays differ in value. Under the FIDE Laws of Chess the side to move may claim a draw when the position is about to occur, or has occurred, for the third time, and the game is drawn automatically at the fifth occurrence; the fifty-move and seventy-five-move rules do the same for the count of quiet moves. A board with a side to move therefore has no value of its own. This is the graph-history interaction problem identified by Campbell~\cite{campbell1985} and given general search-based solutions by Kishimoto and M\"uller~\cite{kishimoto2004,kishimoto2005}.

Engine practice treats the problem heuristically. Stockfish scores a position that has occurred twice before the root as a draw when it is reached again, and scores a position that recurs once inside the search as a draw as well. Tablebase authors have observed that a winning line that follows a rule-aware metric never repeats a position, so that only history supplied from outside the search can spoil a tablebase value~\cite{deman2018}. Neither is a theorem, and neither models the claim a defender may make before executing a move that would repeat a position.

This paper asks a narrow question and answers it exactly. Fix a query consisting of a position, a move-clock value and an arbitrary past, given as the number of earlier occurrences of every position. Must an exact solver from that query keep updating occurrence counts along every hypothetical continuation? It need not. Two frozen sets of positions suffice and are never updated during the solve (Theorem~\ref{thm:main}); under alternating ownership, which chess has, one set suffices after checking the query position's own automatic status (Corollary~\ref{cor:onemask}).

The result is useful in three places. A tablebase probe with a known game history, as in engine analysis of a played game or in adjudication, can be made exact by a solve parametrised by two bitsets rather than by replaying counts. A solved instance carries a certificate, an integer per position, that an independent program checks by local inequalities without search (Corollary~\ref{cor:cert}). And it states precisely why the repetition heuristics in engine practice are value-correct when the game history is known, and how often a history-free probe is wrong (Section~\ref{sec:engine}).

The contributions are as follows.
\begin{enumerate}
\item A fixed-past threshold theorem for a general class of count-triggered neutral draw rules on finite reachability games, and its one-mask specialisation under alternation (Section~\ref{sec:main}).
\item The identification of the latest-winning-clock function with the attractor rank of a reachability game, from which the proof, a key-positional winning strategy, a bucket algorithm linear in the arena size plus the clock bound, a path-budget identity and a search-free certificate follow (Section~\ref{sec:rank}).
\item Machine-checked counterexamples showing that each hypothesis is necessary (Section~\ref{sec:necessity}), and verification of the theorem by two independently written implementations on 850 million queries over exhaustively enumerated arenas (Section~\ref{sec:abstract-exp}).
\item An exact solver for all pawnless three- and four-piece FIDE endgames under the full repetition and move-clock rules, cross-checked against the Syzygy tables, with an exact relation between the deadline and the distance to zeroing (Section~\ref{sec:chess}).
\end{enumerate}
We do not claim priority over the graph-history literature or over the classical fact that a winner can avoid repetition; Section~\ref{sec:related} states what is old and what is new.

\section{Related work}\label{sec:related}

Graph-history interaction was named by Campbell~\cite{campbell1985}; Kishimoto and M\"uller give general solutions for alpha-beta and proof-number search that make stored results reusable across histories~\cite{kishimoto2004,kishimoto2005}, and discuss the DAG method of Schijf, Allis and Uiterwijk in which identical positions are merged after irreversible conversion moves and kept separate elsewhere~\cite{schijf1994}. Thorp and Walden already distinguished the board arrangement from the information needed to continue a history-sensitive game~\cite{thorp1972}. Our result does not replace these search methods; it specialises the rule model and obtains a static parametrisation of a complete solve for one query.

That a winner can avoid repeating a position is classical, going back to Zermelo and K\'alm\'ar~\cite{schwalbe2001}; attractor ranks and positional determinacy of reachability games are standard~\cite{gradel2002,zielonka1998}. In a clock-augmented game an attractor strategy does not repeat an augmented state, which is not the same as not repeating a position; Lemma~\ref{lem:progress} supplies that step. Aminof and Rubin show that memoryless determinacy fails for general cycle objectives~\cite{aminof2017}, which is why the neutrality hypothesis below cannot be dropped.

Clock-sensitive chess tablebases are established, from Haworth's depth-by-the-rule metrics~\cite{haworth2000,haworth2001} and the DTM50 tables of Huntington and Haworth~\cite{huntington2015} to Syzygy, which stores WDL under the fifty-move rule and distance to zeroing~\cite{syzygy}. De Man's observation that DTM50-optimal play avoids new repetitions while supplied history remains a separate concern~\cite{deman2018} is the closest practical antecedent. Our theorem makes the sufficient frozen summary of that supplied past explicit, includes intended-move claims, and states the assumptions at exits; Section~\ref{sec:chess} shows that the deadline is the fifty-move complement of Syzygy's distance to zeroing.

Formally verified endgame databases exist. Hurd verified pawnless four-piece tables in HOL4 without repetition or clock rules~\cite{hurd2005}; Janicic, Maric and Malikovic proved a KRK conjecture by computer~\cite{janicic2019}. Our certificate (Corollary~\ref{cor:cert}) is a ranking certificate in the sense of certifying model checkers~\cite{namjoshi2001} and is checked by an independent program, not by a proof assistant.

Ewerhart's Theorem~2 shows that the value of chess with optional repetition claims, made either at a third occurrence or before the move that would produce one, equals the value of the finite game that ends in a draw at the third occurrence; his appendix defines the finite updateable state, the position together with the set of positions seen once, the set seen at least twice and the quiet count capped at 100, under which subgames are isomorphic~\cite{ewerhart2002}. Both claim forms and the count-set representation are therefore antecedents, and his paper predates the automatic fivefold and seventy-five-move rules. What is new here is the fixed-query reduction, in which, for one query, the seen-once set is discarded, the seen-twice set and the automatic set are frozen instead of updated, the clock enters through a deadline instead of a state coordinate, exits and automatic rules are included, and the result is an integer vector with a linear-time algorithm and a local certificate. Hamkins announces that winning and drawing tactics in chess need no information about repetition or the fifty-move rule~\cite{hamkins2025}; only the introduction of that post is public. Corollary~\ref{cor:positional} is the winner's half of such a statement for arbitrary supplied pasts, with the past entering through the masks; the drawing side is not treated here. Our two-mask statement with prospective claims and exact exits was not found in any source we inspected.

\section{Model}\label{sec:model}

\begin{definition}[Epoch arena]\label{def:arena}
An epoch arena consists of a finite set $Q$ of nonterminal keys, a finite set $T$ of terminals each labelled with an outcome in $\{0,1,\mathrm{D}\}$, an owner $o(q)\in\{0,1\}$ for each key, and for each key two action sets, quiet moves to successors in $Q\cup T$ and exits, each labelled with an exact outcome in $\{0,1,\mathrm{D}\}$. A key with no action is a draw.
\end{definition}

Keys are exact in the sense that owner, actions and terminal status are functions of the key, and the key is the unit that repetition counts. In chess a key is the board, the side to move, the castling entitlements and the legally available en passant capture (Article~9.2.3). Exits abstract irreversible moves into already solved subgames; Proposition~\ref{prop:compose} states when this abstraction is sound.

\begin{definition}[Rules]
Fix clock thresholds $1\le C<H$ and, for every key $x$, repetition thresholds $2\le r_x<a_x$. A query is $(q_0,h_0,c)$ with $q_0\in Q$, $0\le h_0\le H$, and a past count $c:Q\to\N$ excluding the current occurrence of $q_0$.
\end{definition}

\begin{definition}[Explicit-history game]\label{def:explicit}
States are $(q,h,N)$ with $q\in Q\cup T$, $0\le h\le H$ and $N:Q\to\N$ the occurrence counts including the current occurrence; the initial state is $(q_0,h_0,c+\mathbf 1_{q_0})$. A terminal ends the game with its outcome. At a nonterminal, if $h=H$ or $N(q)\ge a_q$ the game ends in a draw; otherwise the owner chooses (a) a quiet move to $q'$, leading to $(q',h+1,N+\mathbf 1_{q'})$, (b) an exit, ending with its outcome, or (c) a claim, ending in a draw, available iff $N(q)\ge r_q$, or some quiet successor $q'$ has $N(q')+1\ge r_{q'}$, or $h\ge C$, or a quiet move exists and $h+1\ge C$. Claims are optional actions.
\end{definition}

Since $h$ increases along quiet moves and every other action terminates, the explicit game is finite and acyclic, and its values are determined by backward induction. Counts may be capped at $a_x$.

\begin{definition}[Frozen-mask game]\label{def:frozen}
For a query let $B=\{x: c(x)\ge r_x-1\}$ and $A=\{x: c(x)\ge a_x-1\}$, so $A\subseteq B$. States are $(q,h)$. At a nonterminal, $h=H$ or $q\in A$ is a draw; a claim is available iff $q\in B$, some quiet successor lies in $B$, $h\ge C$, or a quiet move exists and $h+1\ge C$. Quiet moves lead to $(q',h+1)$. $A$ and $B$ never change.
\end{definition}

The frozen-mask claim set at the initial occurrence of each key is contained in the explicit one; the only omission is a self-loop intended claim with $c(q)=r_q-2$, which never occurs on a winning line (Lemma~\ref{lem:progress}). The frozen-mask game never invents a claim.

\section{Deadlines are attractor ranks}\label{sec:rank}

Fix a target player $P$. $P$ wins a state if $P$ can force a $P$-won terminal or exit; draws and wins for the other player both refute this objective, so $P$ never claims and the opponent may.

\begin{lemma}[Clock monotonicity]\label{lem:mono}
In the frozen-mask game, if $P$ wins $(q,h)$ then $P$ wins $(q,h')$ for every $h'\le h$.
\end{lemma}
\begin{proof}
Couple plays from $(q,h')$ and $(q,h)$ move for move. Quiet moves keep the clocks a constant offset apart; repetition claims depend only on the frozen masks and coincide; clock claims available at the lower clock are available at the higher; exits and terminals coincide. Every opponent action in the lower-clock play is therefore available in the higher-clock play, and the winning strategy transfers.
\end{proof}

Define the deadline $D_P(q)$ as the largest $h$ with $(q,h)$ winning for $P$, or $-1$ if none, with $D_P=H$ at $P$-won terminals and $-1$ at other terminals. Let $\sat(z)=\max(-1,z)$, and for a candidate vector $d$ let the action value be $w_d(e)=\sat(d(q')-1)$ for a quiet move to $q'$, $H-1$ for a $P$-won exit and $-1$ otherwise. Let $K_P(q)=H-1$ at $P$-owned keys and $K_P(q)=\min\{H-1,\,C-1-\mathbf 1[q\text{ has a quiet move}]\}$ at opponent keys, the last clock at which no clock claim is available there. Let $Z_P$ be the set of keys in $A$, actionless keys, and opponent keys at which a repetition claim of Definition~\ref{def:frozen} is available. The operator $F_P$ sets $(F_Pd)(q)=-1$ for $q\in Z_P$ and otherwise
\begin{equation}\label{eq:F}
(F_Pd)(q)=\begin{cases}\min\{K_P(q),\ \max_e w_d(e)\} & o(q)=P,\\[2pt]
\min\{K_P(q),\ \min_e w_d(e)\} & o(q)\ne P,\end{cases}
\end{equation}
and $F_P$ is monotone on the finite lattice $\{-1,\dots,H-1\}^Q$.

\begin{proposition}[Exact deadline equation]\label{prop:fix}
$D_P$ is a fixed point of $F_P$ with the prescribed terminal values, and the only one. Iteration from $-1$ on nonterminals converges to it.
\end{proposition}
\begin{proof}
At a $P$-key, $(q,h)$ is winning iff some action wins at clock $h$; a quiet action wins at $h$ iff $(q',h+1)$ is winning, i.e.\ $h+1\le D_P(q')$; exits win iff $P$-won. At an opponent key every action must win and no claim may be available, which caps $h$ at $K_P(q)$. That is~\eqref{eq:F}. For uniqueness, any fixed point $d$ determines, for each clock $h$, the Boolean layer $\{q: d(q)\ge h\}$, and~\eqref{eq:F} restricted to layers is the backward-induction recursion on the acyclic clock expansion, which has one solution. Monotone iteration on a finite lattice reaches the least fixed point, which is therefore $D_P$.
\end{proof}

\begin{proposition}[Affine identification]\label{prop:rank}
Let $\rho(q)=H-1-D_P(q)$, with $\rho=H$ when $D_P=-1$. Then $\rho$ is the min-max distance to a $P$-winning target on the key arena with unit step cost, node-dependent lower bound $H-1-K_P(q)$ at opponent keys, and target set the $P$-won exits and terminals; that is, $\rho$ is the attractor rank of a reachability game~\cite{gradel2002}, and $(q,h)$ is winning iff $h+\rho(q)\le H-1$.
\end{proposition}
\begin{proof}
Substitute $d=H-1-\rho$ in~\eqref{eq:F}; $\sat(d(q')-1)$ becomes $\min(H,\rho(q')+1)$ and the extrema exchange.
\end{proof}

The floors make $\rho$ a clock-sensitive rank, not the unweighted attractor rank of the key graph. With $H=6$, a defending key whose only action is a $P$-won exit has $D_P=C-1$, so changing $C$ from 2 to 4 changes $\rho$ from 4 to 2 with the graph unchanged.

Standard facts about attractor ranks now yield the rest; Proposition~\ref{prop:bucket} gives the running time. The explicit-history game, by contrast, has up to $(H+1)\prod_x a_x$ states per key.

\begin{lemma}[Strict progress]\label{lem:progress}
At a $P$-owned key with $D_P(q)\ge0$, every action that wins at clock $D_P(q)$ and is quiet satisfies $D_P(q')\ge D_P(q)+1$. At an opponent-owned key with $D_P(q)\ge0$, every quiet move satisfies $D_P(q')\ge D_P(q)+1$.
\end{lemma}
\begin{proof}
Immediate from~\eqref{eq:F}, since a quiet move contributes $\sat(d(q')-1)$ to the extremum that determines $D_P(q)$.
\end{proof}

\begin{corollary}[Key-positional winning strategy]\label{cor:positional}
Let $\sigma_P$ assign to each winning $P$-key an action that wins at clock $D_P(q)$. $\sigma_P$ depends on the key and the masks only; from every winning $(q,h)$ it wins, and along every consistent play the deadline strictly increases, so no key is visited twice.
\end{corollary}
\begin{proof}
The invariant $h\le D_P(q)$ holds initially and is preserved after $\sigma_P$'s quiet move and after any opponent quiet move by Lemma~\ref{lem:progress}; opponent claims are unavailable at keys with $D_P\ge0$ by~\eqref{eq:F} and $h\le K_P(q)$; the play is finite because $h$ increases and $h\le D_P\le H-1$; it ends at a $P$-won terminal or exit because those are the only actions with nonnegative value at the deadline. Strict increase of $D_P$ forbids revisiting a key.
\end{proof}

\begin{remark}\label{rem:current}
An action that wins merely at the current clock need not make progress: a $P$-key with a quiet self-loop and a winning exit can loop while the exit remains available, and looping repeats the key. The deadline-optimal choice is what makes the strategy key-acyclic.
\end{remark}

The deadline can also be read off the caps met along winning continuations, without assuming that it rises by one per move.

\begin{proposition}[Path-clock budget]\label{prop:budget}
Let $\mathcal S(q)$ be the set of key-positional policies for $P$ from $q$ whose continuations, against every defending action, are key-acyclic, avoid $Z_P$ and end only at $P$-won terminals or exits. For $\sigma\in\mathcal S(q)$ let $b(\sigma)=\min_\gamma\min_i\,(K_P(q_i)-i)$ over continuations $\gamma$ and the nonterminal keys $q_i$ on them, where $i$ counts the quiet moves already made from $q$. Then $D_P(q)=\max\{-1,\max_{\sigma\in\mathcal S(q)}b(\sigma)\}$, with $\max\emptyset=-\infty$.
\end{proposition}
\begin{proof}
For $\sigma\in\mathcal S(q)$ the only constraint on a starting clock $h\ge0$ is $h+i\le K_P(q_i)$ at every visited nonterminal, which excludes automatic draws and defending clock claims there; every other outcome is a $P$-win by construction. So every $h\le b(\sigma)$ wins and $D_P(q)\ge b(\sigma)$. Conversely, if $D_P(q)\ge0$ then $\sigma_P$ of Corollary~\ref{cor:positional} lies in $\mathcal S(q)$, and the invariant $h\le D_P(q_i)\le K_P(q_i)$ along every continuation started at $h_0=D_P(q)$ gives $K_P(q_i)-i\ge D_P(q)$, so $b(\sigma_P)\ge D_P(q)$.
\end{proof}

\begin{corollary}[Certificate]\label{cor:cert}
A candidate vector $d\in\{-1,\dots,H-1\}^Q$ with the prescribed terminal values equals $D_P$ iff $d=F_Pd$. Checking this requires only the action sets and~\eqref{eq:F} at each key.
\end{corollary}
\begin{proof}
Proposition~\ref{prop:fix}. $F_P$ maps the lattice into itself, so the range condition is the only side condition.
\end{proof}

Local Boolean consistency of win/loss labels is not a certificate. Two keys $x$ (owned by $P$, with a losing exit) and $y$ (opponent) joined by quiet moves $x\to y\to x$ satisfy ``$x$ has a winning successor and all successors of $y$ are winning'' with both labelled winning, while the true value is a draw. The strict integer decrease in~\eqref{eq:F} excludes such circular support.

\begin{proposition}[Bucket computation]\label{prop:bucket}
Let $n=|Q|$, let $m$ be the number of actions and $T$ the set of terminals. $D_P$ can be computed in $O(n+|T|+m+H)$ bounded-integer operations and $O(n+|T|+m)$ space, and synchronous iteration of~\eqref{eq:F} from the bottom vector reaches $D_P$ after at most $\min(n,H)$ value-changing rounds.
\end{proposition}
\begin{proof}
Compute $\rho=H-1-D_P$ of Proposition~\ref{prop:rank} with buckets indexed $0,\dots,H-1$ and reverse action lists, and write $L(q)=H-1-K_P(q)$ for the floor at $q$. A $P$-won exit or terminal reports cost $0$ to every action entering it; a quiet move to a settled key $x$ reports $\rho(x)+1$; losing actions never report. At a $P$-key each report proposes the candidate $\max(L(q),\text{cost})$; at a defending key a candidate is proposed only when every action has reported, and equals the maximum of the reports and $L(q)$. Candidates of value at least $H$ and keys in $Z_P$ are discarded. Scan the buckets in increasing order; the first removal of a key settles it, later entries for it are ignored, and its reverse edges are processed once. Each action is inspected and reports once, at most $O(m)$ bucket entries are created and the scan costs $O(H)$, which gives the bounds; if $H$ is given in binary the $H$ term is numeric rather than bit-length. For exactness, every proposal is witnessed by a finite winning policy, so no key is settled below its true rank; conversely, by induction on the true rank $k$, the supporting successors of a rank-$k$ key are settled earlier and put a candidate equal to $k$ in bucket $k$ before the scan reaches it. For the round bound, restrict the winning region to the actions of $\sigma_P$ (Corollary~\ref{cor:positional}) and all defending actions. This graph is acyclic by Lemma~\ref{lem:progress}, a branch has at most $n$ keys and at most $H$ distinct nonnegative deadlines, and each synchronous round settles every key whose branches are settled; nonwinning keys stay at $-1$ throughout by monotonicity.
\end{proof}

\section{The fixed-past theorem}\label{sec:main}

\begin{theorem}[Fixed-past threshold equivalence]\label{thm:main}
For every epoch arena, thresholds $(C,H)$ and $(r_x,a_x)_{x\in Q}$, query $(q_0,h_0,c)$ and target $P$: $P$ wins the explicit-history game from the query iff $h_0\le D_P(q_0)$ in the frozen-mask game with $A=\{c\ge a-1\}$ and $B=\{c\ge r-1\}$. Applying this to both players, the explicit win/draw/loss value equals the frozen-mask value.
\end{theorem}
\begin{proof}
($\Leftarrow$) Play $\sigma_P$ of Corollary~\ref{cor:positional} in the explicit game. Along any consistent play no key repeats, so at the current key $q$ the explicit count is exactly $c(q)+1$; a current claim or automatic draw exists in the explicit game only if it exists in the frozen-mask game, where it does not at keys with $D_P\ge0$. A defender's intended claim concerns a quiet successor $q'$. If $q'$ were on the play so far, strict progress along the play gives $D_P(q')\le D_P(q)$, contradicting Lemma~\ref{lem:progress}; so $q'$ is fresh, its count after the move would be $c(q')+1$, and an intended claim exists only if $q'\in B$, which~\eqref{eq:F} already excludes at winning defending keys. Clock claims coincide. Hence every explicit action available to the opponent has been accounted for, and the play ends in a $P$-win.

($\Rightarrow$) If $h_0>D_P(q_0)$ the opponent has a counterstrategy in the frozen-mask game, which is finite. Every frozen-mask claim is a legal explicit claim, because explicit counts are at least the frozen counts plus the current occurrence; explicit automatic draws only help the opponent; moves and exits coincide. Projecting the explicit play to $(q,h)$, the counterstrategy is well defined and prevents a $P$-win.
\end{proof}

\begin{corollary}[Value equivalence of pasts]
Two pasts inducing the same $(A,B)$ give every query with the same key and clock the same value, even when individual counts differ.
\end{corollary}

\begin{corollary}[One mask under alternation]\label{cor:onemask}
If every quiet move between keys changes the owner, then for $q_0\notin A$ the deadlines computed with $A=\emptyset$ agree with $D_P$ at every key outside $A$; the query needs only $B$ after checking $q_0\notin A$.
\end{corollary}
\begin{proof}
A winning strategy cannot enter a defending key in $A\subseteq B$, since a current claim exists there, nor a $P$-key in $A$, since by alternation its predecessor is defending and has an intended claim into $B$.
\end{proof}

\begin{proposition}[Rank-separated composition]\label{prop:compose}
Let a finite game have exact keys with an integer rank $\rho$ that quiet moves preserve and reset moves strictly decrease, terminals and legal actions functions of the key, and supplied histories that are valid plays. Then the game can be solved by increasing rank, treating each reset successor as a fresh zero-clock query in the already solved lower-rank arena, and Theorem~\ref{thm:main} applies within each rank.
\end{proposition}
\begin{proof}
After a reset every future key has strictly lower rank than every key before it, so no future key equals a discarded past key or a key of the accrued prefix, and no intended repetition claim can concern a reset successor; the reset also resets the clock, so no intended clock claim concerns it. By induction on rank the lower-rank outcomes are exact, so resets are exact exits in the sense of Definition~\ref{def:arena}.
\end{proof}

The no-return condition is thus a property of the abstraction from a full game to epoch arenas, discharged by Proposition~\ref{prop:compose}, not a hypothesis of Theorem~\ref{thm:main}. Without it the abstraction is unsound: an exact value for the state after a reset does not prevent the opponent from claiming, before the reset is executed, by an intended move back into the prefix. The released code contains an explicit instance.

\begin{remark}
The masks preserve a value query, not the transition system. Two pasts with equal masks diverge after one move, and a certificate for one mask pair says nothing about another. A live application must keep enough history to form fresh masks for each new query.
\end{remark}

\section{Necessity of the hypotheses}\label{sec:necessity}

Each example was checked by computing both games explicitly; the scripts are released.

\begin{example}[Neutrality, rewarded repetition]
One $P$-key with a quiet self-loop and no exit, $r=2$, empty past. If the second occurrence wins for the mover, the explicit value is a $P$-win; the frozen-mask game, in which repetition draws, gives a draw.
\end{example}

\begin{example}[Neutrality, penalised repetition]
Keys $x$ ($P$) $\to y$ (opponent) $\to x$, no exits, $r=3$, $c(x)=1$. If creating the third occurrence loses, the opponent is forced back to $x$ and loses; with $B=\emptyset$ the frozen-mask game sees a bare cycle and returns a draw.
\end{example}

\begin{example}[Clock monotonicity]
$x$ ($P$, with a self-loop and an edge to $y$), $y$ (opponent) $\to z$ ($P$) $\to$ winning exit, $H=6$, and the opponent may claim only when $h\in\{1,2\}$. Explicitly the query is a draw, since moving to $y$ early allows a claim and waiting at $x$ triggers the automatic repetition draw. A frozen-mask analysis with monotone clock claims reports a win. Waiting is profitable only when a larger clock can help, which monotonicity forbids.
\end{example}

\begin{example}[Alternation, for Corollary~\ref{cor:onemask}]
$x\to y$ with both keys owned by $P$, a winning exit at $y$, and $y\in A$. With the automatic mask the arrival at $y$ is a draw; without it $P$ wins. No defending turn intervenes. In the enumeration of Section~\ref{sec:abstract-exp} the one-mask reduction fails on exactly the 12{,}288 non-alternating cases and on none of the 354{,}902 alternating ones.
\end{example}

\section{Verification of the reduction}\label{sec:abstract-exp}

Two implementations were written from Definitions~\ref{def:explicit} and~\ref{def:frozen}. The first computes the explicit-history value by memoised backward induction over $(q,h,N)$; the second iterates~\eqref{eq:F}; a third computes the frozen-mask value by backward induction over $(q,h)$ without deadlines. The two-key family has, per key, 2 owners, 4 quiet-successor subsets and 8 exit subsets, giving 64 nonterminal configurations, plus 3 terminals; the three-key family restricts exits to at most one outcome. Every arena is queried at every key, every clock and every past-count vector in $\{0,\dots,a-1\}^{|Q|}$, including vectors no legal play produces.

\begin{table}[t]\centering\footnotesize
\caption{Explicit-history versus frozen-mask values. ``Proof-scenario hits'' counts winning queries in which, along the key-positional strategy, some defending key has a quiet successor lying in the supplied past one occurrence below the claim threshold and off the realised play, the case the forward direction of Theorem~\ref{thm:main} turns on; with $r=2$ it cannot occur.}
\label{tab:abstract}
\begin{tabular}{@{}llrrr@{}}\toprule
Family & $(r,a,C,H)$ & Queries & Hits & Mism.\\ \midrule
2-key exhaustive & (2,3,2,3) & 323{,}208 & n/a & 0\\
2-key exhaustive & (3,5,4,6) & 1{,}571{,}150 & n/a & 0\\
random, 1--9 keys & varied & 295{,}265 & n/a & 0\\
dense random, 3--6 keys & varied & 1{,}257{,}754 & n/a & 0\\
3-key exhaustive & (2,3,2,3) & 93{,}083{,}904 & 0 & 0\\
3-key exhaustive & (3,5,4,6) & 754{,}152{,}000 & 30{,}720 & 0\\ \bottomrule
\end{tabular}
\end{table}

Table~\ref{tab:abstract} reports zero disagreements on 850 million queries. This is consistency between two implementations of the same definitions and not a proof; it also says nothing about whether a particular game's rules are modelled correctly, which Section~\ref{sec:chess} addresses separately.

The implementations also let a count sit on a terminal slot, in which case a quiet move to that terminal at threshold carries an intended claim in both games; queries with such counts are included in the totals and never disagreed. The propositions of Sections~\ref{sec:rank} and~\ref{sec:chess} that go beyond the theorem were executed as well, with a fixed seed. The bucket computation of Proposition~\ref{prop:bucket} was compared with synchronous iteration of~\eqref{eq:F} on every two-key arena and mask of both families and on 20{,}000 random arenas with up to 12 keys and random thresholds, 345{,}252 deadline vectors in all, with no difference. For the certificate of Corollary~\ref{cor:cert}, the true vector passed the local test in every one of 40{,}000 cases and every one of 2{,}486{,}548 single-entry mutations of it was rejected. The budget identity of Proposition~\ref{prop:budget} was checked by enumerating every key-positional policy on 3{,}000 arenas with up to six keys, 17{,}884 queries, with no difference from the deadline. For Proposition~\ref{prop:lift}, 20{,}000 random encoded arenas with up to three representatives per key and duplicated actions gave the projected deadlines in every case. The arrival game of Proposition~\ref{prop:arrival} agreed with the explicit-history value on 258{,}455 queries over 10{,}000 random alternating arenas. The script and its output are released with the paper.

\section{FIDE chess}\label{sec:chess}

\subsection{Rule mapping}
The normative rules are the FIDE Laws effective 1 January 2023~\cite{fide}. Table~\ref{tab:rules} gives the mapping. Halfmove counts are plies since the last capture or pawn move. A position is the same as an earlier one when the same player is to move, the pieces stand on the same squares, and the same moves are available, including castling entitlements and an en passant capture that is actually legal (Articles 9.2.2 and 9.2.3). Losing a castling right changes the key but does not reset the clock.

\begin{table}[t]\centering\footnotesize
\caption{FIDE rules and model parameters.}
\label{tab:rules}
\begin{tabular}{@{}lll@{}}\toprule
Rule & Model & Article\\ \midrule
Same position & key exactness & 9.2.2, 9.2.3\\
Threefold claim (current, intended) & $r=3$ & 9.2.1\\
Fifty-move claim (current, intended) & $C=100$ & 9.3\\
Fivefold draw & $a=5$ & 9.6.1\\
Seventy-five-move draw & $H=150$ & 9.6.2\\
Capture, pawn move & reset & 3.1, 3.7\\
Mate, stalemate, dead position & terminals & 5.1.1, 5.2\\ \bottomrule
\end{tabular}
\end{table}

For the rank of Proposition~\ref{prop:compose} let $N(q)$ be the number of pieces and set $\rho(q)=N(q)+\sum_{\text{white pawns}}(8-j)+\sum_{\text{black pawns}}(j-1)$ over pawn ranks $j$. Every capture or pawn move strictly decreases $\rho$, including en passant and promotion, which removes a positive pawn term even when the piece count is unchanged; every other move preserves it. Since chess alternates the side to move, Corollary~\ref{cor:onemask} applies: after checking whether the query position has already occurred four times, the only historical parameter an exact solve needs is the set of positions that have occurred at least twice.

An incomplete but sound dead-position recogniser does not change values: a set of keys closed under legal successors from which neither side can be mated has value draw whether or not the game is stopped on entry, so declaring only a subset of dead positions terminal is safe for win/draw/loss. It does not reproduce the moment of termination, which the solver below does not need.

\subsection{Solver}
A solver implementing~\eqref{eq:F} directly was written for pawnless materials with up to four pieces. Keys are the piece squares and the side to move. En passant cannot arise without pawns, but castling entitlements can, since a king on e1 with an unmoved rook on h1 is pawnless; the solver covers positions without castling rights, so a key that still carries an entitlement is outside its tables, and it is a different key from the same board without the entitlement. Quiet moves are non-captures; captures are exits whose outcome is the fresh zero-clock value in the lower-material table, in the order of Proposition~\ref{prop:compose}. Terminals are checkmate and stalemate. Masks are supplied as key lists; the fresh case has empty masks. Iteration is synchronous over all keys until no value changes.

Keys are stored with labelled pieces, so a position with two identical pieces has two encodings, and repetition counts physical positions. The following makes labelled storage sound.

\begin{proposition}[Labelled encodings]\label{prop:lift}
Let $\pi\colon\tilde Q\cup\tilde T\to Q\cup T$ be a surjection from an encoded arena onto the physical one that preserves ownership, terminal status and outcomes, and such that the projected quiet successors and exits of every encoded key are exactly those of its physical key, with multiplicity allowed. With $\tilde A=\pi^{-1}(A)$ and $\tilde B=\pi^{-1}(B)$, the encoded deadlines are $\tilde D_P=D_P\circ\pi$.
\end{proposition}
\begin{proof}
For every vector $d$ on $Q$, the extrema in~\eqref{eq:F} at $\tilde q$ over the encoded actions equal those at $\pi(\tilde q)$, since duplicated actions change neither a minimum nor a maximum; caps, terminal values and membership in $Z_P$ are preserved by construction. Hence $\tilde F_P(d\circ\pi)=(F_Pd)\circ\pi$, and uniqueness of the fixed point (Proposition~\ref{prop:fix}) gives the claim.
\end{proof}

The masks must therefore be formed from physical positions and lifted to every encoding; a mask that lists one encoding of a repeated position is wrong. Table~\ref{tab:chess} counts labelled keys, so for KQQvK, KRRvK, KBBvK and KNNvK the number of physical positions is half the number shown.

Table~\ref{tab:chess} reports, for every pawnless three- and four-piece material, the number of legal keys, the win/draw/loss counts for the side to move, the number of iterations and the time on a desktop CPU.

\begin{table*}[t]\centering\small
\caption{Pawnless materials solved under the FIDE repetition and move-clock rules with fresh history. Win, draw and loss are for White to move; legal keys count both sides to move. Rounds are iterations of~\eqref{eq:F} until no value changes; time is on an eight-core desktop.}
\label{tab:chess}
\begin{tabular}{@{}lrrrrrr@{}}\toprule
Material & Legal keys & Win & Draw & Loss & Rounds & s\\ \midrule
KQvK & 368{,}452 & 144{,}508 & 0 & 0 & 17 & 1\\
KRvK & 399{,}112 & 175{,}168 & 0 & 0 & 26 & 2\\
KBvK & 417{,}228 & 0 & 193{,}284 & 0 & 1 & 0\\
KNvK & 429{,}440 & 0 & 205{,}496 & 0 & 1 & 0\\
KQvKQ & 17{,}905{,}216 & 3{,}737{,}092 & 5{,}174{,}888 & 40{,}628 & 14 & 54\\
KQvKR & 19{,}733{,}336 & 8{,}863{,}768 & 71{,}704 & 17{,}136 & 45 & 142\\
KQvKB & 20{,}785{,}072 & 8{,}925{,}252 & 27{,}356 & 0 & 22 & 66\\
KQvKN & 21{,}487{,}864 & 8{,}894{,}128 & 58{,}480 & 0 & 30 & 63\\
KRvKR & 21{,}561{,}456 & 3{,}139{,}232 & 7{,}569{,}032 & 72{,}464 & 8 & 24\\
KRvKB & 22{,}613{,}192 & 3{,}787{,}160 & 6{,}993{,}568 & 0 & 23 & 59\\
KRvKN & 23{,}315{,}984 & 5{,}210{,}920 & 5{,}569{,}800 & 8 & 34 & 63\\
KBvKB & 23{,}664{,}928 & 416 & 11{,}831{,}936 & 112 & 2 & 7\\
KBvKN & 24{,}367{,}720 & 16 & 11{,}832{,}440 & 8 & 2 & 6\\
KNvKN & 25{,}070{,}512 & 40 & 12{,}535{,}208 & 8 & 2 & 4\\
KQQvK & 19{,}317{,}704 & 5{,}657{,}120 & 0 & 0 & 8 & 16\\
KQRvK & 20{,}571{,}880 & 6{,}911{,}296 & 0 & 0 & 9 & 18\\
KQBvK & 21{,}359{,}016 & 7{,}698{,}432 & 0 & 0 & 13 & 28\\
KQNvK & 21{,}905{,}880 & 8{,}245{,}296 & 0 & 0 & 15 & 29\\
KRRvK & 21{,}985{,}768 & 8{,}325{,}184 & 0 & 0 & 10 & 20\\
KRBvK & 23{,}027{,}424 & 9{,}366{,}840 & 0 & 0 & 21 & 42\\
KRNvK & 23{,}565{,}632 & 9{,}905{,}048 & 0 & 0 & 21 & 32\\
KBBvK & 23{,}824{,}640 & 5{,}007{,}216 & 5{,}156{,}840 & 0 & 30 & 57\\
KBNvK & 24{,}536{,}088 & 10{,}822{,}184 & 53{,}320 & 0 & 47 & 78\\
KNNvK & 25{,}159{,}888 & 1{,}232 & 11{,}498{,}072 & 0 & 2 & 3\\
\bottomrule
\end{tabular}
\end{table*}

\subsection{Cross-check against Syzygy}
The Syzygy tables store, for each position, whether the side to move wins, draws or loses under the fifty-move rule, with wins that need more than fifty moves marked as cursed, and the distance to the next zeroing move~\cite{syzygy}. In our model a win that requires more than one hundred quiet plies is a draw, because the defender claims at the hundredth ply, so a cursed win maps to a draw. Every legal position of every three-piece material, and a uniform random sample of 300{,}000 positions of every four-piece material, was probed through python-chess and compared with the fresh-history value of our tables.

\begin{table*}[t]\centering\small
\caption{Agreement with Syzygy on fresh-history win/draw/loss (three-piece materials in full, four-piece materials on a uniform random sample), and the value of $D+\mathrm{DTZ}$ on winning positions: 100 when Syzygy's rounded DTZ50'' equals the exact plies to zeroing, 99 when it is one less, 150 when the winning move captures or mates immediately.}
\label{tab:syzygy}
\begin{tabular}{@{}lrrrrr@{}}\toprule
Material & Compared & Disagree & $D{+}\mathrm{DTZ}{=}100$ & ${=}99$ & ${=}150$\\ \midrule
KQvK & 368{,}452 & 0 & 142{,}060 & 0 & 2{,}448\\
KRvK & 399{,}112 & 0 & 173{,}656 & 0 & 1{,}512\\
KBvK & 417{,}228 & 0 & 0 & 0 & 0\\
KNvK & 429{,}440 & 0 & 0 & 0 & 0\\
KQvKQ & 300{,}000 & 0 & 6{,}622 & 0 & 118{,}903\\
KQvKR & 300{,}000 & 0 & 80{,}271 & 1{,}377 & 100{,}057\\
KQvKB & 300{,}000 & 0 & 76{,}913 & 0 & 51{,}494\\
KQvKN & 300{,}000 & 0 & 74{,}520 & 0 & 49{,}737\\
KRvKR & 300{,}000 & 0 & 3{,}514 & 0 & 84{,}235\\
KRvKB & 300{,}000 & 0 & 0 & 10{,}176 & 39{,}776\\
KRvKN & 300{,}000 & 0 & 0 & 28{,}200 & 38{,}760\\
KBvKB & 300{,}000 & 0 & 0 & 0 & 6\\
KBvKN & 300{,}000 & 0 & 0 & 0 & 2\\
KNvKN & 300{,}000 & 0 & 0 & 0 & 1\\
KQQvK & 297{,}902 & 0 & 39{,}090 & 25{,}279 & 23{,}835\\
KQRvK & 298{,}583 & 0 & 34{,}545 & 53{,}130 & 12{,}942\\
KQBvK & 299{,}424 & 0 & 57{,}140 & 44{,}941 & 5{,}690\\
KQNvK & 300{,}000 & 0 & 76{,}877 & 31{,}942 & 3{,}954\\
KRRvK & 298{,}706 & 0 & 40{,}823 & 63{,}598 & 9{,}117\\
KRBvK & 300{,}000 & 0 & 99{,}810 & 20{,}898 & 1{,}217\\
KRNvK & 300{,}000 & 0 & 65{,}486 & 59{,}645 & 1{,}095\\
KBBvK & 299{,}487 & 0 & 0 & 62{,}440 & 159\\
KBNvK & 300{,}000 & 0 & 0 & 132{,}137 & 28\\
KNNvK & 300{,}000 & 0 & 0 & 0 & 7\\
\midrule
Total & 7{,}608{,}334 & 0 & & & \\
\bottomrule
\end{tabular}
\end{table*}

Agreement on win, draw and loss is exact (Table~\ref{tab:syzygy}). The deadline also determines the distance to zeroing. From a winning position with the winner to move, let $t$ be the number of plies to the first capture or mate when both sides play deadline-optimally, the winner choosing a move that wins at the deadline and the defender the move that leaves the winner the smallest deadline. If the winner's first move captures or mates, $t=1$ and $D=149$, because a terminal or exit win is not subject to the clock. Otherwise the defender's cap $K_P=98$ applies at the first defending key, every later ply raises the deadline by exactly one, and $D=100-t$. If instead the defender's only legal moves are captures, the cap is $K_P=99$ and the relation is unchanged. The deadline does not rise by one per ply throughout, as Proposition~\ref{prop:budget} anticipates: with White Kc6 and Rb1 against Black Ka8, White to move, the forced line 1.\,Kc7 Ka7 2.\,Ra1 mate has deadlines $97$, $98$ and $149$ at its three nonterminal positions, the last because a mating move is not subject to the clock. A playout over 3{,}000 sampled winning positions of every four-piece material, or all of them where fewer exist, found $D+t\in\{100,150\}$ and nothing else, in 52{,}200 positions. Syzygy's DTZ50'' is documented as equal to the exact distance or one less because of rounding in the tables~\cite{syzygy}; accordingly $D+\mathrm{DTZ}$ takes only the values 100, 99 and 150 in Table~\ref{tab:syzygy}. The deadline is thus the fifty-move complement of the exact distance to zeroing, and it is the quantity that composes with masks.

\subsection{Engine practice, explained and bounded}\label{sec:engine}
Stockfish scores a position as a draw when it is reached in the search and has occurred twice before in the game, and when it recurs once strictly after the root. Theorem~\ref{thm:main} explains why this is value-correct when the game history is available, and where it stops being so.

The arrival rule is the intended-move clause in disguise. If a position $x$ has occurred twice and the search reaches it by a move of the defender, the exact game lets the defender claim a draw before making that move (Article~9.2.1), and the frozen-mask game does the same through the intended claim at the predecessor; scoring $x$ as a draw on arrival gives the predecessor the same value. If $x$ is reached by a move of the winner, the defender to move at $x$ holds a current claim, and the exact value at $x$ is a draw as well. In both cases arrival scoring equals the frozen-mask value with $B=\{x\}$. Scoring a second occurrence inside the search as a draw affects only lines in which some player repeats a position after the root; by Corollary~\ref{cor:positional} the winner never needs to, so the exact winning value survives, and the defender gains nothing from a line the winner declines. Within one epoch and with the full game history, the heuristic and the theorem therefore agree on win/draw/loss. The precise statement follows.

\begin{proposition}[Arrival scoring]\label{prop:arrival}
Assume alternation. Let the arrival game be the frozen-mask game with these changes: at the root, automatic status is checked and a current claim is retained when $q_0\in B$; a nonterminal key reached by a quiet move at positive depth is drawn on arrival when it lies in $B$ or has already occurred in the current line; no other repetition claim exists. The arrival game has the root win/draw/loss value of the explicit-history game.
\end{proposition}
\begin{proof}
By Theorem~\ref{thm:main} it suffices to match the frozen-mask game. $\sigma_P$ never repeats a key and, as in Corollary~\ref{cor:onemask}, never enters a key of $B$ at positive depth, so the added arrival draws never stop it and $P$ wins the arrival game whenever $h_0\le D_P(q_0)$. Conversely a frozen-mask counterstrategy transfers, because a current claim at a positive-depth defending key in $B$ is an arrival draw, an intended claim into $B$ is an arrival draw one ply later, root claims are retained, and every other arrival draw only prevents a $P$-win. Both directions hold for either target.
\end{proof}

Drawing at the root whenever $q_0\in B$ would be wrong, since the owner of a root with a winning exit declines its claim.

They part company when the history is absent. A probe from a bare FEN uses $B=\emptyset$. Take White Kc6 and Rb1 against Black Ka8 with Black to move; the fresh table gives $D_W=94$, a win. If the game has already cycled 1.\,\ldots Ka7 2.\,Rb2 Ka8 3.\,Rb1 twice with Black declining to claim, the position has occurred twice before, Black holds a current claim, and the exact value is a draw, $D_W=-1$; the same board at the same clock has two different values, and the bare probe reports the wrong one. Whenever a winning position's every winning continuation passes through a defending key with a quiet successor that has occurred twice, the exact value is a draw and the history-free value is a win. The frequency of this dependence is measurable from the tables. Call a won position with the winner to move fragile if it has exactly one winning move at clock zero; a single earlier occurrence of the resulting position, which the winner's opponent may then claim, turns it into a draw, and no alternative exists. Over 1{,}602{,}200 sampled won positions of the four-piece materials, 44{,}257 (2.8\%) are fragile. The fraction is zero in every material where one side has two pieces against a bare king, below one percent in KQvKR, KQvKB and KQvKN, 3.7\% in KQvKQ, 4.0\% in KRvKR, and 17.7\% in KRvKB and 17.3\% in KRvKN, where the rook side's win often passes through a single quiet move. Mating moves and winning captures are not counted, since a terminal or a lower-material position cannot have occurred before. Fragility is a property of the fresh table, so a probe with history can be answered exactly by the mask, but a probe without history is wrong on every fragile position whose successor the game has already visited twice.

\section{Limitations}
The theorem preserves win/draw/loss for reachability objectives with neutral draw rules and monotone clocks. It does not preserve shortest-mate distances, does not apply to repetition rules that reward or penalise a player, such as superko or perpetual-check sanctions, and requires exact exits, which for chess means solved lower-material tables. The chess solver covers pawnless materials up to four pieces; positions with pawns need castling and en passant in the key and the pawn terms in the rank, which the model supports but the solver does not implement. Verification is by agreement between separately written programs and against Syzygy, not by a proof assistant. The algorithms are explicit-arena algorithms; a compactly specified game with an enormous arena remains enormous to enumerate, and nothing here bounds the number of distinct mask pairs that a sequence of queries can produce or gives their values a small shared representation.

\section{Conclusion}
For finite games with count-triggered neutral draw rules and a monotone clock, the exact value of a query with an arbitrary past depends on that past only through two frozen threshold sets, one under alternation. The proof is an attractor-rank argument on positions, which also gives a positional winning strategy and a search-free certificate. The reduction holds on 850 million exhaustively generated queries, and its chess instantiation reproduces the Syzygy tables on every compared pawnless position, with the deadline equal to the fifty-move complement of the distance to zeroing.

\section*{Reproducibility}
The abstract solvers, the counterexample scripts, the property-check script for the propositions, the chess solver, the solved tables' summaries and the Syzygy comparison outputs are released with the paper.

\begin{thebibliography}{99}\small
\bibitem{campbell1985} M. Campbell, ``The graph-history interaction: on ignoring position history,'' in \emph{Proc. ACM Annual Conf.}, 1985, pp. 278--280.
\bibitem{kishimoto2004} A. Kishimoto and M. M\"uller, ``A general solution to the graph history interaction problem,'' in \emph{Proc. AAAI}, 2004, pp. 644--649.
\bibitem{kishimoto2005} A. Kishimoto and M. M\"uller, ``A solution to the GHI problem for depth-first proof-number search,'' \emph{Inf. Sci.}, vol. 175, no. 4, pp. 296--314, 2005.
\bibitem{schijf1994} M. Schijf, L. V. Allis, and J. W. H. M. Uiterwijk, ``Proof-number search and transpositions,'' \emph{ICCA J.}, vol. 17, no. 2, pp. 63--74, 1994.
\bibitem{thorp1972} E. Thorp and W. E. Walden, ``A computer assisted study of Go on M$\times$N boards,'' \emph{Inf. Sci.}, vol. 4, no. 1, pp. 1--33, 1972.
\bibitem{schwalbe2001} U. Schwalbe and P. Walker, ``Zermelo and the early history of game theory,'' \emph{Games Econ. Behav.}, vol. 34, pp. 123--137, 2001.
\bibitem{gradel2002} E. Gr\"adel, W. Thomas, and T. Wilke, Eds., \emph{Automata, Logics, and Infinite Games}, LNCS 2500. Springer, 2002.
\bibitem{zielonka1998} W. Zielonka, ``Infinite games on finitely coloured graphs with applications to automata on infinite trees,'' \emph{Theor. Comput. Sci.}, vol. 200, pp. 135--183, 1998.
\bibitem{aminof2017} B. Aminof and S. Rubin, ``First-cycle games,'' \emph{Inf. Comput.}, vol. 254, pp. 195--216, 2017.
\bibitem{haworth2000} G. McC. Haworth, ``Strategies for constrained optimisation,'' \emph{ICGA J.}, vol. 23, no. 1, pp. 9--20, 2000.
\bibitem{haworth2001} G. McC. Haworth, ``Depth by the rule,'' \emph{ICGA J.}, vol. 24, no. 3, p. 160, 2001.
\bibitem{huntington2015} G. Huntington and G. McC. Haworth, ``Depth to mate and the 50-move rule,'' \emph{ICGA J.}, vol. 38, no. 2, 2015.
\bibitem{syzygy} R. de Man, ``Syzygy endgame tablebases,'' software and documentation, \url{https://github.com/syzygy1/tb}.
\bibitem{deman2018} R. de Man, reply in ``Syzygy question,'' TalkChess, 22 May 2018, \url{https://talkchess.com/viewtopic.php?p=762997}.
\bibitem{hurd2005} J. Hurd, ``Formal verification of chess endgame databases,'' in \emph{Theorem Proving in Higher Order Logics: Emerging Trends}, Oxford Univ. Computing Lab. Tech. Rep. PRG-RR-05-02, 2005, pp. 85--100.
\bibitem{janicic2019} P. Jani\v{c}i\'c, F. Mari\'c, and M. Malikovi\'c, ``Computer-assisted proving of combinatorial conjectures over finite domains: A case study of a chess conjecture,'' \emph{Log. Methods Comput. Sci.}, vol. 15, no. 1, 2019.
\bibitem{namjoshi2001} K. S. Namjoshi, ``Certifying model checkers,'' in \emph{Proc. CAV}, LNCS 2102, 2001, pp. 2--13.
\bibitem{hamkins2025} J. D. Hamkins, ``Tactics versus strategies: the case of chess,'' \emph{Infinitely More}, 17 Aug. 2025, introduction public, remainder behind a subscription.
\bibitem{ewerhart2002} C. Ewerhart, ``Backward induction and the game-theoretic analysis of chess,'' \emph{Games Econ. Behav.}, vol. 39, no. 2, pp. 206--214, 2002, doi: 10.1006/game.2001.0900.
\bibitem{fide} FIDE, \emph{Laws of Chess}, effective 1 Jan. 2023, \url{https://handbook.fide.com/chapter/E012023}.
\end{thebibliography}
\end{document}
